% Root and DE independent full proof audits passed on 2026-11-15.
\begin{statement}{C.38}{Proved}{Unbounded prime-power arithmetic-trivial characters}
Let \(Q\) be the full gentle group modulo the closed normal
subgroup generated by its absolute Galois image.
For every prime \(p\geq5\), every integer \(k\geq1\), and every
even integer \(m\geq12\) divisible by \((p-1)p^{k-1}\), there is
a continuous surjection
\[
                  Q^{\mathrm{ab}}\longrightarrow
                       (\mathbb Z/p^k\mathbb Z)^{d_m},
                                                               \tag{C38.1}
\]
where \(d_m=a_m-b_m>0\) is the dimension lower bound of
C.34--C.36.  The independent values are witnessed on global
multiplier one.  The integers \(d_m\) tend to infinity along
the indicated progression.

Consequently, for every \(p\geq5\),
\[
 \dim_{\mathbb F_p}\operatorname{Hom}_{\mathrm{cts}}
                         (Q,\mathbb F_p)=\aleph_0,
                                                               \tag{C38.2}
\]
and there is a simultaneous continuous surjection
\[
 Q^{\mathrm{ab}}\longrightarrow
       \left(\prod_{j\geq1}C_2\right)
       \times
       \prod_{\substack{p\geq5\\p\ {\rm prime}}}
             \prod_{k,j\geq1}\mathbb Z/p^k\mathbb Z .
                                                               \tag{C38.3}
\]
In particular its \(p\)-primary quotients have unbounded
exponent for every \(p\geq5\).  Neither a
\(\mathbb Z_p\) quotient nor the complete abelianization is
asserted.  The general characters in this statement are not
asserted to vanish on the pentagon group.
\end{statement}
\begin{proof}
Fix \(p\geq5\) and a finite top degree \(m\).
Write \(\mathcal G\) for the actual local gentle image,
\(\mathcal U\) for its local multiplier kernel, and
\(\mathcal U_n\) for the normalized conjugator filtration.
The complete integral lifting argument of C.33 applies
in every finite degree, and gives
\[
       \mathcal U_n/\mathcal U_{n+1}=S_n(\mathbb Z_p).
                                                               \tag{C38.4}
\]
These are the actual full-group images, not just solutions
of the equations modulo a nilpotent term.
Set
\[
       U=\mathcal U/\mathcal U_{m+1},\qquad
       G=\mathcal G/\mathcal U_{m+1}.
\]
The action on the integral Magnus algebra modulo
augmentation degree \(m+2\) realizes \(U\) faithfully as
a compact group of unipotent matrices over \(\mathbb Z_p\).
For \(u\in U\), the operator \(u-1\) raises degree by at
least three.  This follows from \(S_2=0\), as in C.33.
No bound on the nilpotency class relative to \(p\) is imposed.

\paragraph{The actual rational Lie algebra.}
The finite matrix logarithm and exponential are inverse
polynomials over \(\mathbb Q_p\).  Define
\[
             L=\operatorname{span}_{\mathbb Q_p}
                               \{\log u:u\in U\}.
                                                               \tag{C38.5}
\]
This is a Lie subalgebra of the finite-dimensional nilpotent
matrix algebra.  Here is a direct verification that avoids
an integral correspondence in class at least \(p\).
For \(X=\log u\) and \(Y=\log v\), the finite polynomial
\(\log(u^a v^b)=\operatorname{BCH}(aX,bY)\) belongs to \(L\)
whenever \(a,b\) are nonnegative integers.  Evaluating on a
sufficiently large finite integer grid and inverting the
two rational Vandermonde matrices shows that every
coefficient of this polynomial belongs to \(L\).
Its coefficient of \(ab\) is \([X,Y]/2\).
Thus \([X,Y]\in L\); bilinearity proves the assertion.

The group \(U\) contains a neighborhood of the identity
in \(\exp L\).  Choose a basis \(X_1,\ldots,X_d\) of \(L\)
consisting of logarithms of actual elements
\(u_1,\ldots,u_d\in U\).
For every \(t_i\in\mathbb Z_p\), the continuous power
\(u_i^{t_i}\) is in \(U\), and equals \(\exp(t_iX_i)\):
the equality holds on the dense subset of integers and
both sides are continuous.
In these basis coordinates the polynomial map
\[
 (t_1,\ldots,t_d)\longmapsto
       \log\bigl(u_1^{t_1}\cdots u_d^{t_d}\bigr)
                                                               \tag{C38.6}
\]
has the form \(t+R(t)\), where \(R\) has no term of
degree zero or one.
For sufficiently large \(r\), the map \(R\) sends
\(p^r\mathbb Z_p^d\) into \(p^{r+1}\mathbb Z_p^d\) and is
Lipschitz there with constant at most \(p^{-1}\).
For each \(y\) in that ball, the contraction
\(t\mapsto y-R(t)\) therefore gives a solution of
\(t+R(t)=y\).
It follows that
\[
       \exp\left(p^r\sum_i\mathbb Z_pX_i\right)\subseteq U .
                                                               \tag{C38.7}
\]
This is an explicit rational argument; no claim that
\(\log U\) itself is an integral Lie lattice is needed.

Let \(L_n\) be the intersection of \(L\) with the
operators that raise augmentation degree by at least \(n\).
This is the same filtration as the normalized conjugator
filtration: a first degree-\(n\) conjugator term \(f_n\)
changes \(Y\) first by \([Y,f_n]\), and the centralizer
of \(Y\) in the characteristic-zero free Lie algebra
has no homogeneous part in degree \(n\geq2\).
The leading terms of logarithm and exponential agree
with their linear terms.
Equation (C38.7) shows that every vector of \(L_n\),
after multiplying by a sufficiently large power of \(p\),
is the logarithm of an element of \(U_n\).
Conversely logarithms of \(U_n\) lie in \(L_n\).
Taking leading terms and using (C38.4) therefore gives
\[
           L_n/L_{n+1}=S_n(\mathbb Q_p).
                                                               \tag{C38.8}
\]
The induced bracket is the actual leading Ihara bracket.
This verifies the rational Lie filtration directly from
the completed group image.

\paragraph{A grading supplied by an actual wild multiplier.}
Choose an element \(g_0\) in the actual local Galois image
whose multiplier is \(1+p\), and put
\(A=\operatorname{Ad}(g_0)|_L\).
It preserves the filtration and acts on its degree-\(n\)
quotient by the scalar \((1+p)^n\).
These scalars, for \(3\leq n\leq m\), are pairwise
distinct in \(\mathbb Q_p\).
The filtration therefore implies
\[
           \prod_{n=3}^{m}\bigl(A-(1+p)^n\bigr)=0.
                                                               \tag{C38.9}
\]
For example, apply the factors in ascending filtration
order; each annihilates the current successive quotient
and moves the image to the next term.
The polynomial has no repeated root, so \(A\) is
semisimple and yields a rational splitting
\[
      L=\bigoplus_{n=3}^{m}L^{(n)},\qquad
      L^{(n)}\xrightarrow{\ \sim\ }S_n(\mathbb Q_p).
                                                               \tag{C38.10}
\]
Since \(A\) is a Lie automorphism, the bracket of the
\(i\) and \(j\) eigenspaces lies in the \(i+j\)
eigenspace when \(i+j\leq m\), and is zero otherwise.
Its leading term identifies it with the Ihara bracket
under (C38.10).
Thus the entire rational Lie algebra has the gentle
grading.  This removes hidden higher filtered bracket
relations even when many tame weights coincide and the
nilpotency class exceeds \(p\).

The weight-\(m\) part of \(L/[L,L]\) has dimension
\[
        h_m\geq a_m-b_m=d_m.
                                                               \tag{C38.11}
\]
Indeed the number of independent pair brackets in that
degree is at most \(b_m\), exactly as in C.34.
Choose \(h_m\) independent rational linear functionals
on this part and extend them by zero on the other
eigenspaces.  Together they give a surjection
\(\ell:L\to\mathbb Q_p^{h_m}\) annihilating \([L,L]\).
Consequently
\[
                      u\longmapsto\ell(\log u)
                                                               \tag{C38.12}
\]
is an ordinary continuous group homomorphism:
every non-linear term of the finite BCH polynomial
is a Lie commutator and is annihilated.

Conjugation on \(L/[L,L]\) factors through
\(G/U=\mathbb Z_p^\times\).
Indeed inner conjugation by \(u=\exp X\) induces
\(\exp(\operatorname{ad}X)\), which is the identity
on the Lie abelianization.
These conjugation operators commute with \(A\) and
therefore preserve its distinct weight spaces.
On the top weight \(m\), they act exactly as on
the top filtration quotient, namely by
\(\lambda_p(g)^m\).
Thus
\[
       \ell(\operatorname{Ad}(g)X)
                    =\lambda_p(g)^m\ell(X)
             \quad(g\in G,\ X\in L).
                                                               \tag{C38.13}
\]

\paragraph{Integral values and finite prime-power characters.}
The image of (C38.12) is a compact additive subgroup
\(J\subset\mathbb Q_p^{h_m}\).
It is a \(\mathbb Z_p\)-module: it is closed and stable
under integer multiplication.
It spans \(\mathbb Q_p^{h_m}\), since \(\log U\)
spans \(L\) and \(\ell\) is onto.
It is consequently a full free \(\mathbb Z_p\)-lattice.
Choose a basis of \(J\).  In this basis (C38.12) becomes
a continuous surjection
\[
                \Phi:U\longrightarrow\mathbb Z_p^{h_m},
       \qquad
       \Phi(gug^{-1})=\lambda_p(g)^m\Phi(u).
                                                               \tag{C38.14}
\]
This passage uses the actual compact image, so no
unsupported integral BCH assertion or reduction of a
possibly nonprimitive rational functional is involved.

Suppose \((p-1)p^{k-1}\mid m\).
For every \(\lambda\in\mathbb Z_p^\times\),
\(\lambda^m\equiv1\pmod{p^k}\).
Reduction of (C38.14) is therefore an invariant
surjection \(U\to(\mathbb Z/p^k\mathbb Z)^{h_m}\).
The continuous five-term sequence extends each coordinate
to \(G\), because
\[
 H^2_{\mathrm{cts}}(\mathbb Z_p^\times,\mathbb Z/p^k\mathbb Z)=0.
                                                               \tag{C38.15}
\]
As before, this follows from averaging on \(\mu_{p-1}\)
and the \(p\)-cohomological dimension one of
\(1+p\mathbb Z_p\).

The restriction of each extended character to
\(G_{\mathbb Q}\) is unramified away \(p\), by the
actual pro-\(p\) tangential action cited in C.33.
Kronecker--Weber identifies every such character with
a scalar multiple of the cyclotomic character
\[
 c_{p,k}(\lambda)=
        \frac{\log(\lambda_p^{p-1})}{p}\pmod{p^k}.
                                                               \tag{C38.16}
\]
Here the logarithm is taken on \(1+p\mathbb Z_p\).
It is an ordinary additive character and is onto:
the map \(\log/p:1+p\mathbb Z_p\to\mathbb Z_p\)
is an isomorphism, while \(p-1\) is a unit.
Subtracting these scalar multiples removes the
arithmetic restrictions simultaneously.
The correction is zero on global multiplier one.
Every element of the actual local multiplier-one image
has a full global multiplier-one lift by P.4.
Thus the joint surjectivity in (C38.14), reduced
modulo \(p^k\), remains valid for the corrected
full-group characters.
This proves (C38.1), since \(h_m\geq d_m\).
C.34 gives \(d_m\to\infty\) along the progression,
and C.36 gives positivity for every even \(m\geq12\).

\paragraph{Choosing all finite orders simultaneously.}
Let \(K\) be the global multiplier-one subgroup.
For fixed \(p,k\), let \(V_{p,k}\) be the subspace of
\(\operatorname{Hom}_{\mathrm{cts}}(K,\mathbb F_p)\)
consisting of reductions of the restrictions of
arithmetic-trivial full-group characters to
\(\mathbb Z/p^k\mathbb Z\).
The joint surjections just proved show that
\(\dim_{\mathbb F_p}V_{p,k}=\infty\).
Enumerate all pairs \((k,j)\) of positive integers.
At each step choose an actual such character whose
mod-\(p\) restriction lies outside the span of all
previously chosen restrictions.
This is possible because the relevant \(V_{p,k}\)
is infinite-dimensional and only finitely many
characters have been chosen.

For any finite collection of the resulting characters,
the image \(H\) in its product target \(T\) satisfies
\(H+pT=T\), by independence of their mod-\(p\)
restrictions to \(K\).
The finite \(\mathbb Z_p\)-module \(T/H\) therefore
equals \(p(T/H)\), so it is zero.
The finite product map is onto already on \(K\).
Compactness now gives a surjection onto
\(\prod_{k,j\geq1}\mathbb Z/p^k\mathbb Z\).
Combine these maps for distinct primes and C.30's
independent quadratic family: every finite target
product has pairwise coprime primary factors and is
onto.  Another use of compactness proves (C38.3).

Finally the full gentle group is a closed subgroup
of the automorphism group of a finitely generated
free profinite group, hence is countably based.
It has only countably many clopen subsets, and thus
only countably many continuous maps to any fixed
finite group.  This gives the upper bound
\(\aleph_0\) in (C38.2); the preceding construction
gives the matching lower bound.
The product of finite cyclic groups in (C38.3)
is not being replaced by a \(\mathbb Z_p\) factor:
unbounded finite orders alone do not provide a
compatible infinite tower of characters.
\end{proof}
